% ============================================================
% PHỤ LỤC TỔNG HỢP KIẾN THỨC CHƯƠNG 1--4
% Học phần: Lý thuyết điều khiển tự động
%
% File độc lập để \input vào cuối main.tex sau này.
% KHÔNG chứa \documentclass, \begin{document}, \end{document}.
% KHÔNG sửa main.tex.
% ============================================================

\providecommand{\Lap}{\mathcal{L}}
\providecommand{\jj}{\mathrm{j}}

\clearpage
\begin{center}
{\LARGE\bfseries PHỤ LỤC TỔNG HỢP KIẾN THỨC}\\[0.25em]
{\Large\bfseries CHƯƠNG 1 -- CHƯƠNG 4}\\[0.45em]
{\small Lý thuyết, bảng công thức, quy trình giải bài tập và các lỗi thường gặp}
\end{center}
\hrule
\vspace{0.7em}

\noindent\textbf{Quy ước dùng trong phụ lục.}
Trong các công thức Laplace, giả sử tín hiệu bằng $0$ với $t<0$.
Khi xác định hàm truyền, điều kiện đầu được lấy bằng $0$.
Ký hiệu $r(t)$: tín hiệu đặt; $e(t)$: sai lệch; $u(t)$: tín hiệu điều khiển;
$c(t)$ hoặc $y(t)$: tín hiệu ra; $R(s),E(s),U(s),C(s),Y(s)$ là các ảnh Laplace tương ứng.

% ============================================================
\clearpage
\section*{CHƯƠNG 1 -- TỔNG QUAN VỀ HỆ THỐNG ĐIỀU KHIỂN TỰ ĐỘNG}

\subsection*{1. Khái niệm cốt lõi}

\textbf{Điều khiển} là quá trình thu thập thông tin, xử lý thông tin và tác động lên hệ thống
để đáp ứng thực tế gần với mục tiêu đã định trước.

\textbf{Điều khiển tự động} là quá trình điều khiển trong đó thao tác điều chỉnh được thực hiện
không cần sự can thiệp trực tiếp liên tục của con người.

\medskip
\noindent Một hệ điều khiển điển hình có thể nhận diện qua chuỗi:
\[
\boxed{\text{tín hiệu đặt}}
\longrightarrow
\boxed{\text{so sánh}}
\longrightarrow
\boxed{\text{bộ điều khiển}}
\longrightarrow
\boxed{\text{đối tượng}}
\longrightarrow
\boxed{\text{tín hiệu ra}} .
\]

\subsection*{2. Các tín hiệu thường gặp}

\begin{center}
\small
\begin{tabular}{|p{0.18\textwidth}|p{0.25\textwidth}|p{0.46\textwidth}|}
\hline
\textbf{Ký hiệu} & \textbf{Tên} & \textbf{Ý nghĩa}\\
\hline
$r(t)$ & tín hiệu đặt / tín hiệu chuẩn & giá trị mong muốn của đại lượng cần điều khiển\\
\hline
$e(t)$ & sai lệch & chênh lệch giữa giá trị đặt và tín hiệu hồi tiếp\\
\hline
$u(t)$ & tín hiệu điều khiển & đầu ra bộ điều khiển, tác động lên đối tượng\\
\hline
$c(t)$ hoặc $y(t)$ & tín hiệu ra & đại lượng cần điều khiển / quan sát\\
\hline
$c_{ht}(t)$ & tín hiệu hồi tiếp & tín hiệu đo từ đầu ra đưa trở về bộ so sánh\\
\hline
$d(t)$ & nhiễu & tác động ngoài không mong muốn lên hệ thống\\
\hline
\end{tabular}
\end{center}

Với hồi tiếp âm:
\[
\boxed{e(t)=r(t)-c_{ht}(t)}.
\]
Nếu cảm biến / khâu hồi tiếp là đơn vị thì $c_{ht}(t)=c(t)$:
\[
\boxed{e(t)=r(t)-c(t)}.
\]

\subsection*{3. Hệ hở và hệ kín}

\textbf{Hệ hở:}
đầu ra không được dùng trực tiếp để điều chỉnh lại tín hiệu điều khiển.
Ưu điểm là cấu trúc đơn giản; nhược điểm là khả năng tự sửa sai và khử nhiễu kém.

\textbf{Hệ kín:}
đầu ra được đo, hồi tiếp và so sánh với tín hiệu đặt.
Sai lệch được dùng để sinh tín hiệu điều khiển mới.

\medskip
\noindent Ba nguyên tắc điều khiển thường gặp:
\begin{itemize}[leftmargin=1.7em,itemsep=0.2em]
  \item \textbf{Điều khiển bù nhiễu:} đo hoặc ước lượng nhiễu rồi bù tác động của nhiễu.
  \item \textbf{Điều khiển theo sai lệch:} so sánh đầu ra với tín hiệu đặt rồi điều chỉnh để giảm sai lệch.
  \item \textbf{Điều khiển phối hợp:} kết hợp bù nhiễu và phản hồi sai lệch.
\end{itemize}

\subsection*{4. Phân loại hệ thống -- bảng nhớ nhanh}

\begin{center}
\small
\begin{tabular}{|p{0.30\textwidth}|p{0.59\textwidth}|}
\hline
\textbf{Tiêu chí} & \textbf{Nhóm thường gặp}\\
\hline
Theo số ngõ vào -- ngõ ra & SISO: một vào -- một ra; MIMO: nhiều vào -- nhiều ra\\
\hline
Theo dạng tín hiệu & hệ liên tục; hệ rời rạc / lấy mẫu / số\\
\hline
Theo tính tuyến tính & hệ tuyến tính; hệ phi tuyến\\
\hline
Theo sự thay đổi theo thời gian & hệ bất biến theo thời gian; hệ biến thiên theo thời gian\\
\hline
Theo cấu trúc phản hồi & hệ hở; hệ kín\\
\hline
Theo mục tiêu điều khiển & ổn định hóa; theo chương trình; bám / theo dõi; thích nghi; tối ưu\\
\hline
\end{tabular}
\end{center}

\subsection*{5. Dạng bài Chương 1 và cách làm}

\subsubsection*{Dạng 1: Xác định các phần tử của một hệ điều khiển}

Làm theo thứ tự:
\begin{enumerate}[leftmargin=1.8em,itemsep=0.2em]
  \item Xác định \textbf{đại lượng cần giữ / cần bám} $\Rightarrow$ đầu ra $c(t)$.
  \item Xác định \textbf{giá trị mong muốn} $\Rightarrow$ tín hiệu đặt $r(t)$.
  \item Xác định \textbf{đối tượng} cần điều khiển.
  \item Xác định \textbf{cơ cấu chấp hành} tạo tác động vật lý lên đối tượng.
  \item Xác định \textbf{cảm biến / khâu đo}.
  \item Xác định \textbf{bộ điều khiển}.
  \item Chỉ ra \textbf{nhiễu} nếu đề có.
\end{enumerate}

\subsubsection*{Dạng 2: Nhận biết hệ hở hay hệ kín}

Tìm xem có đường đưa thông tin đầu ra quay trở lại bộ so sánh hay không.
Có phản hồi $\Rightarrow$ hệ kín; không có phản hồi $\Rightarrow$ hệ hở.

\subsubsection*{Dạng 3: Vẽ sơ đồ khối từ mô tả vật lý}

Không vẽ ngay chi tiết thiết bị. Trước hết rút hệ thống về logic:
\[
r(t)\rightarrow \text{bộ so sánh}\rightarrow \text{bộ điều khiển}
\rightarrow \text{đối tượng}\rightarrow c(t),
\]
sau đó mới thêm cảm biến, phản hồi và nhiễu.

\medskip
\noindent\textbf{Lỗi thường gặp:}
nhầm \emph{đối tượng điều khiển} với \emph{cơ cấu chấp hành};
bỏ sót dấu $-$ tại bộ so sánh của hồi tiếp âm.

% ============================================================
\clearpage
\section*{CHƯƠNG 2 -- MÔ HÌNH TOÁN HỌC, LAPLACE, HÀM TRUYỀN VÀ KHÔNG GIAN TRẠNG THÁI}

\subsection*{1. Từ hệ vật lý đến phương trình vi phân}

Quan hệ vào -- ra của hệ tuyến tính bất biến liên tục thường được viết:
\[
a_n\frac{d^ny}{dt^n}
+a_{n-1}\frac{d^{n-1}y}{dt^{n-1}}
+\cdots+a_1\frac{dy}{dt}+a_0y
=
b_m\frac{d^mu}{dt^m}
+\cdots+b_1\frac{du}{dt}+b_0u.
\]
Thông thường với hệ hợp thức: $n\ge m$.

\subsection*{2. Biến đổi Laplace}

Định nghĩa:
\[
\boxed{
F(s)=\Lap\{f(t)\}=\int_0^\infty f(t)e^{-st}\,dt,
\qquad s=\sigma+\jj\omega.
}
\]

\subsection*{3. Bảng Laplace cơ bản}

\begin{center}
\small
\renewcommand{\arraystretch}{1.35}
\begin{tabular}{|p{0.40\textwidth}|p{0.42\textwidth}|}
\hline
\textbf{Hàm thời gian $f(t)$} & \textbf{Ảnh Laplace $F(s)$}\\
\hline
$\delta(t)$ & $1$\\
\hline
$1(t)$ & $\dfrac{1}{s}$\\
\hline
$t\,1(t)$ & $\dfrac{1}{s^2}$\\
\hline
$t^n1(t)$ & $\dfrac{n!}{s^{n+1}}$\\
\hline
$e^{-at}1(t)$ & $\dfrac{1}{s+a}$\\
\hline
$\sin \omega t\,1(t)$ & $\dfrac{\omega}{s^2+\omega^2}$\\
\hline
$\cos \omega t\,1(t)$ & $\dfrac{s}{s^2+\omega^2}$\\
\hline
$e^{-at}\sin\omega t\,1(t)$ & $\dfrac{\omega}{(s+a)^2+\omega^2}$\\
\hline
$e^{-at}\cos\omega t\,1(t)$ & $\dfrac{s+a}{(s+a)^2+\omega^2}$\\
\hline
$\dfrac{t^{n-1}}{(n-1)!}e^{-at}1(t)$ & $\dfrac{1}{(s+a)^n}$\\
\hline
\end{tabular}
\renewcommand{\arraystretch}{1}
\end{center}

\subsection*{4. Các tính chất Laplace cần thuộc}

Nếu $F(s)=\Lap\{f(t)\}$ và $G(s)=\Lap\{g(t)\}$:

\textbf{Tuyến tính}
\[
\Lap\{af(t)+bg(t)\}=aF(s)+bG(s).
\]

\textbf{Dịch thời gian}
\[
\Lap\{f(t-T)1(t-T)\}=e^{-Ts}F(s),\qquad T>0.
\]

\textbf{Dịch theo biến $s$}
\[
\Lap\{e^{-at}f(t)\}=F(s+a).
\]

\textbf{Nhân với $t^n$}
\[
\boxed{\Lap\{t^nf(t)\}=(-1)^n\frac{d^nF(s)}{ds^n}}.
\]

\textbf{Ảnh của đạo hàm}
\[
\Lap\{f'(t)\}=sF(s)-f(0^+),
\]
\[
\Lap\{f''(t)\}=s^2F(s)-sf(0^+)-f'(0^+),
\]
và tổng quát:
\[
\boxed{
\Lap\{f^{(n)}(t)\}
=s^nF(s)-s^{n-1}f(0^+)-s^{n-2}f'(0^+)-\cdots-f^{(n-1)}(0^+).
}
\]

\textbf{Ảnh của tích phân}
\[
\Lap\left\{\int_0^t f(\tau)\,d\tau\right\}=\frac{F(s)}{s}.
\]

\textbf{Định lý giá trị đầu}
\[
f(0^+)=\lim_{s\to\infty}sF(s),
\]
khi các giả thiết hội tụ thỏa mãn.

\textbf{Định lý giá trị cuối}
\[
\boxed{f(\infty)=\lim_{s\to0}sF(s)}
\]
chỉ dùng khi các cực liên quan của $sF(s)$ nằm trong nửa mặt phẳng trái
(với các ngoại lệ biên phải được xét riêng).

\subsection*{5. Laplace ngược bằng phân tích phân thức}

Cho
\[
F(s)=\frac{B(s)}{A(s)}.
\]
Nếu $\deg B\ge\deg A$, phải chia đa thức trước để đưa về phân thức hợp thức.

\subsubsection*{Trường hợp cực thực đơn}

Nếu
\[
F(s)=\sum_{i=1}^{n}\frac{C_i}{s-p_i},
\]
thì
\[
\boxed{C_i=\left[(s-p_i)F(s)\right]_{s=p_i}},
\qquad
f(t)=\sum_i C_ie^{p_it}1(t).
\]

\subsubsection*{Trường hợp cực bội}

Nếu $p$ là cực bội $r$:
\[
F(s)=\frac{A_r}{(s-p)^r}+\frac{A_{r-1}}{(s-p)^{r-1}}
+\cdots+\frac{A_1}{s-p}+\cdots
\]
Đặt
\[
H(s)=(s-p)^rF(s).
\]
Khi đó:
\[
A_r=H(p),\qquad
A_{r-1}=H'(p),\qquad
A_{r-2}=\frac{H''(p)}{2!},\ldots
\]
Tổng quát:
\[
\boxed{
A_k=
\frac{1}{(r-k)!}
\left[
\frac{d^{\,r-k}}{ds^{\,r-k}}H(s)
\right]_{s=p},
\qquad k=1,\ldots,r.
}
\]

\subsubsection*{Trường hợp cặp cực phức}

Đưa mẫu về dạng:
\[
(s+a)^2+\omega^2.
\]
Tử nên tách:
\[
A(s+a)+B\omega.
\]
Sau đó dùng:
\[
\Lap^{-1}\left\{\frac{s+a}{(s+a)^2+\omega^2}\right\}
=e^{-at}\cos\omega t,
\]
\[
\Lap^{-1}\left\{\frac{\omega}{(s+a)^2+\omega^2}\right\}
=e^{-at}\sin\omega t.
\]

\subsection*{6. Quy trình giải phương trình vi phân bằng Laplace}

\begin{enumerate}[leftmargin=1.9em,itemsep=0.2em]
  \item Lấy Laplace hai vế.
  \item Thay công thức đạo hàm, giữ đúng các điều kiện đầu.
  \item Gom các hạng chứa $Y(s)$ và rút $Y(s)$.
  \item Phân tích $Y(s)$ thành các phân thức tối giản.
  \item Lấy Laplace ngược.
  \item Kiểm tra điều kiện đầu và thế lại phương trình nếu cần.
\end{enumerate}

\textbf{Ví dụ nhanh}
\[
y'(t)+4y(t)=5\,1(t),\qquad y(0)=0.
\]
Lấy Laplace:
\[
sY(s)+4Y(s)=\frac{5}{s}.
\]
Do đó:
\[
Y(s)=\frac{5}{s(s+4)}
=\frac54\frac1s-\frac54\frac1{s+4}.
\]
Suy ra:
\[
\boxed{y(t)=\frac54\left(1-e^{-4t}\right)1(t)}.
\]

\subsection*{7. Hàm truyền đạt}

Với hệ LTI, khi các điều kiện đầu bằng $0$:
\[
\boxed{G(s)=\frac{Y(s)}{U(s)}}.
\]

Nếu
\[
a_ny^{(n)}+\cdots+a_1\dot y+a_0y
=
b_mu^{(m)}+\cdots+b_1\dot u+b_0u,
\]
thì:
\[
\boxed{
G(s)=
\frac{b_ms^m+b_{m-1}s^{m-1}+\cdots+b_1s+b_0}
{a_ns^n+a_{n-1}s^{n-1}+\cdots+a_1s+a_0}.
}
\]

\subsubsection*{Quy trình từ phương trình vi phân sang hàm truyền}

\begin{enumerate}[leftmargin=1.9em,itemsep=0.2em]
  \item Xác định rõ đầu vào và đầu ra.
  \item Cho tất cả điều kiện đầu bằng $0$.
  \item Lấy Laplace hai vế.
  \item Gom $Y(s)$ một phía, $U(s)$ phía còn lại.
  \item Lấy tỷ số $G(s)=Y(s)/U(s)$ và rút gọn.
\end{enumerate}

\textbf{Ví dụ}
\[
\ddot y+3\dot y+2y=u
\quad\Longrightarrow\quad
(s^2+3s+2)Y(s)=U(s)
\]
nên:
\[
\boxed{G(s)=\frac{1}{s^2+3s+2}}.
\]

\subsection*{8. Mô hình các phần tử vật lý cơ bản}

\subsubsection*{Mạch điện}

\begin{center}
\small
\begin{tabular}{|p{0.17\textwidth}|p{0.35\textwidth}|p{0.29\textwidth}|}
\hline
\textbf{Phần tử} & \textbf{Miền thời gian} & \textbf{Quan hệ miền $s$, điều kiện đầu $0$}\\
\hline
Điện trở $R$ & $v_R(t)=Ri_R(t)$ & $V_R(s)=RI_R(s)$, $Z_R=R$\\
\hline
Cuộn cảm $L$ & $v_L(t)=L\dfrac{di_L}{dt}$ & $V_L(s)=sL\,I_L(s)$, $Z_L=sL$\\
\hline
Tụ điện $C$ & $i_C(t)=C\dfrac{dv_C}{dt}$ & $I_C(s)=sC\,V_C(s)$, $Z_C=\dfrac1{sC}$\\
\hline
\end{tabular}
\end{center}

Với mạch RLC nối tiếp, lấy điện áp trên tụ làm đầu ra:
\[
LC\ddot u_C+RC\dot u_C+u_C=u_{in},
\]
nên:
\[
\boxed{
\frac{U_C(s)}{U_{in}(s)}
=
\frac{1}{LCs^2+RCs+1}.
}
\]

\subsubsection*{Hệ cơ tịnh tiến}

\begin{center}
\small
\begin{tabular}{|p{0.18\textwidth}|p{0.35\textwidth}|p{0.28\textwidth}|}
\hline
\textbf{Phần tử} & \textbf{Quan hệ lực -- chuyển vị} & \textbf{Miền $s$}\\
\hline
Khối lượng $M$ & $F_M=M\ddot x$ & $F_M=Ms^2X$\\
\hline
Giảm chấn $B$ & $F_B=B\dot x$ & $F_B=BsX$\\
\hline
Lò xo $K$ & $F_K=Kx$ & $F_K=KX$\\
\hline
\end{tabular}
\end{center}

Với hệ khối lượng -- giảm chấn -- lò xo:
\[
M\ddot x+B\dot x+Kx=F(t)
\]
nên:
\[
\boxed{\frac{X(s)}{F(s)}=\frac{1}{Ms^2+Bs+K}}.
\]

\subsection*{9. Không gian trạng thái}

Dạng chuẩn:
\[
\boxed{
\dot x=Ax+Bu,\qquad y=Cx+Du.
}
\]

Từ mô hình trạng thái sang hàm truyền:
\[
\boxed{G(s)=C(sI-A)^{-1}B+D}.
\]

Phương trình đặc trưng của ma trận trạng thái:
\[
\boxed{\det(sI-A)=0}.
\]
Các nghiệm là các trị riêng của $A$; với một biểu diễn tối thiểu, chúng trùng với các cực của hàm truyền.

\subsubsection*{Quy trình chuyển hàm truyền sang trạng thái}

Một cách thường dùng:
\begin{enumerate}[leftmargin=1.9em,itemsep=0.2em]
  \item Đưa hàm truyền về phương trình vi phân.
  \item Chọn các biến trạng thái, thường chọn theo chuỗi đạo hàm của đầu ra.
  \item Viết các phương trình vi phân cấp một.
  \item Sắp xếp vào dạng ma trận $\dot x=Ax+Bu$, $y=Cx+Du$.
\end{enumerate}

\medskip
\noindent\textbf{Lưu ý:}
một hàm truyền có thể có nhiều mô hình trạng thái tương đương; mô hình trạng thái không duy nhất.

% ============================================================
\clearpage
\section*{CHƯƠNG 3 -- ĐẠI SỐ SƠ ĐỒ KHỐI VÀ GRAPH TÍN HIỆU}

\subsection*{1. Ba phần tử của sơ đồ khối}

\begin{itemize}[leftmargin=1.7em,itemsep=0.2em]
  \item \textbf{Khối chức năng:} nếu đầu vào là $X(s)$, đầu ra là $Y(s)$ thì
  \[
  Y(s)=G(s)X(s).
  \]
  \item \textbf{Bộ tổng:} đầu ra bằng tổng đại số các tín hiệu vào theo dấu $+$ / $-$.
  \item \textbf{Điểm rẽ nhánh:} các nhánh lấy ra từ cùng một điểm có cùng giá trị tín hiệu.
\end{itemize}

\subsection*{2. Các cấu trúc phải thuộc}

\textbf{Nối tiếp}
\[
\boxed{G_{eq}(s)=G_1(s)G_2(s)\cdots G_n(s)}.
\]

\textbf{Song song}
\[
\boxed{G_{eq}(s)=G_1(s)+G_2(s)+\cdots+G_n(s)}
\]
có xét dấu đúng tại bộ tổng.

\textbf{Hồi tiếp âm}
\[
\boxed{
G_{cl}(s)=\frac{G(s)}{1+G(s)H(s)}.
}
\]

\textbf{Hồi tiếp dương}
\[
\boxed{
G_{cl}(s)=\frac{G(s)}{1-G(s)H(s)}.
}
\]

Hồi tiếp đơn vị: đặt $H(s)=1$.

\subsection*{3. Quy tắc di chuyển bộ tổng và điểm rẽ nhánh}

Giả sử có một khối $G(s)$.

\begin{center}
\small
\renewcommand{\arraystretch}{1.35}
\begin{tabular}{|p{0.38\textwidth}|p{0.45\textwidth}|}
\hline
\textbf{Phép di chuyển} & \textbf{Điều chỉnh nhánh phụ để giữ tương đương}\\
\hline
Bộ tổng: từ trước $G$ ra sau $G$ & nhánh phụ phải nhân $G$\\
\hline
Bộ tổng: từ sau $G$ về trước $G$ & nhánh phụ phải chia $G$\\
\hline
Điểm rẽ: từ trước $G$ ra sau $G$ & nhánh lấy ra phải qua $1/G$\\
\hline
Điểm rẽ: từ sau $G$ về trước $G$ & nhánh lấy ra phải qua $G$\\
\hline
\end{tabular}
\renewcommand{\arraystretch}{1}
\end{center}

\subsection*{4. Quy trình rút gọn sơ đồ khối nhiều vòng}

\begin{enumerate}[leftmargin=1.9em,itemsep=0.2em]
  \item Đánh dấu các vòng hồi tiếp \textbf{nhỏ nhất / trong cùng} trước.
  \item Gộp các khối nối tiếp hoặc song song nhìn thấy ngay.
  \item Nếu chưa gộp được, di chuyển bộ tổng hoặc điểm rẽ để tạo cấu trúc chuẩn.
  \item Rút gọn vòng trong, sau đó tiến dần ra vòng ngoài.
  \item Sau mỗi phép biến đổi, kiểm tra lại vị trí bộ tổng, dấu và nhánh lấy tín hiệu.
\end{enumerate}

\subsubsection*{Mẹo kiểm tra nhanh}

Hai sơ đồ tương đương khi quan hệ vào -- ra của chúng không đổi.
Nếu nghi ngờ một phép di chuyển, đặt một tín hiệu ký hiệu $X$ tại nút trước biến đổi,
tính tín hiệu tại các nút trước và sau; hai kết quả phải trùng nhau.

\subsection*{5. Graph tín hiệu -- khái niệm}

\begin{itemize}[leftmargin=1.7em,itemsep=0.2em]
  \item \textbf{Nút}: đại diện cho một tín hiệu.
  \item \textbf{Nhánh có hướng}: mang hệ số truyền của nhánh.
  \item \textbf{Đường truyền thuận $P_k$}: đường từ nút vào đến nút ra, không đi qua một nút quá một lần.
  \item \textbf{Vòng lặp $L_i$}: đường khép kín bắt đầu và kết thúc tại cùng một nút.
  \item \textbf{Hai vòng không chạm nhau}: không có nút chung.
\end{itemize}

\subsection*{6. Công thức Mason}

\[
\boxed{
\frac{C(s)}{R(s)}
=
\frac{\displaystyle\sum_k P_k\Delta_k}{\Delta}.
}
\]

Trong đó:
\[
\boxed{
\Delta
=
1-\sum_iL_i
+\sum_{\substack{i<j\\\text{không chạm}}}L_iL_j
-\sum_{\substack{i<j<m\\\text{đôi một không chạm}}}L_iL_jL_m
+\cdots
}
\]

$\Delta_k$ được tính giống $\Delta$ nhưng \textbf{loại tất cả vòng chạm đường truyền thuận $P_k$}.

\subsection*{7. Quy trình giải graph bằng Mason}

\begin{enumerate}[leftmargin=1.9em,itemsep=0.2em]
  \item Liệt kê tất cả đường truyền thuận $P_1,P_2,\ldots$.
  \item Liệt kê tất cả vòng $L_1,L_2,\ldots$, ghi cả dấu của từng vòng.
  \item Tìm các cặp, bộ ba, \ldots vòng không chạm nhau.
  \item Tính $\Delta$.
  \item Với từng $P_k$, loại vòng chạm $P_k$ để tính $\Delta_k$.
  \item Thay vào công thức Mason và rút gọn.
\end{enumerate}

\medskip
\noindent\textbf{Lỗi thường gặp:}
bỏ dấu âm của nhánh hồi tiếp; nhầm khái niệm hai vòng không chạm nhau với hai vòng không dùng chung nhánh
(trong Mason, không chạm nghĩa là \textbf{không có nút chung});
bỏ sót đường truyền thuận; dùng một vòng chạm $P_k$ trong $\Delta_k$.

% ============================================================
\clearpage
\section*{CHƯƠNG 4 -- ĐẶC TÍNH ĐỘNG HỌC VÀ CÁC KHÂU ĐỘNG HỌC CƠ BẢN}

\subsection*{1. Đặc tính thời gian}

Cho hệ có hàm truyền $G(s)$, điều kiện đầu bằng $0$.

\textbf{Đáp ứng xung / hàm trọng lượng}
\[
r(t)=\delta(t)\Longrightarrow R(s)=1,
\]
\[
C(s)=G(s),
\qquad
\boxed{g(t)=\Lap^{-1}\{G(s)\}}.
\]

\textbf{Đáp ứng bước / hàm quá độ}
\[
r(t)=1(t)\Longrightarrow R(s)=\frac1s,
\]
\[
C(s)=\frac{G(s)}s,
\qquad
\boxed{h(t)=\Lap^{-1}\left\{\frac{G(s)}s\right\}}.
\]

Quan hệ:
\[
\boxed{
h(t)=\int_0^t g(\tau)\,d\tau,
\qquad
g(t)=\frac{dh(t)}{dt}.
}
\]

Do đó, nếu biết $h(t)$:
\[
H(s)=\Lap\{h(t)\},
\qquad
\boxed{G(s)=sH(s)}
\]
khi $h(0^-)=0$.

\subsection*{2. Đặc tính tần số}

Thay:
\[
s=\jj\omega
\]
vào hàm truyền:
\[
\boxed{
G(\jj\omega)
=
P(\omega)+\jj Q(\omega)
=
M(\omega)e^{\jj\varphi(\omega)}.
}
\]

Biên độ:
\[
\boxed{
M(\omega)=|G(\jj\omega)|
=\sqrt{P^2(\omega)+Q^2(\omega)}.
}
\]

Pha:
\[
\boxed{
\varphi(\omega)=\arg G(\jj\omega).
}
\]

Biên độ theo dB:
\[
\boxed{
L(\omega)=20\log_{10}M(\omega)\quad\text{dB}.
}
\]

Nếu:
\[
G(s)=K\prod_iG_i(s),
\]
thì:
\[
L(\omega)=20\log_{10}|K|+\sum_iL_i(\omega),
\]
\[
\varphi(\omega)=\arg K+\sum_i\varphi_i(\omega).
\]

\subsection*{3. Bảng các khâu động học điển hình}

\begin{center}
\scriptsize
\renewcommand{\arraystretch}{1.4}
\begin{tabular}{|p{0.19\textwidth}|p{0.22\textwidth}|p{0.27\textwidth}|p{0.20\textwidth}|}
\hline
\textbf{Khâu} & \textbf{$G(s)$} & \textbf{Đáp ứng bước chính} & \textbf{Đặc tính tần số chính}\\
\hline
Khuếch đại / tỉ lệ & $K$ & $c(t)=K\,1(t)$ & $M=|K|$, pha $0^\circ$ nếu $K>0$\\
\hline
Tích phân lý tưởng & $\dfrac1s$ & $h(t)=t\,1(t)$ & $M=1/\omega$, pha $-90^\circ$\\
\hline
Vi phân lý tưởng & $s$ & $h(t)=\delta(t)$ & $M=\omega$, pha $+90^\circ$\\
\hline
Quán tính bậc nhất & $\dfrac{K}{Ts+1}$ & $h(t)=K(1-e^{-t/T})1(t)$ & góc pha $0^\circ\to-90^\circ$\\
\hline
Vi phân bậc nhất & $K(Ts+1)$ & có thành phần xung tại $t=0$ & góc pha $0^\circ\to+90^\circ$\\
\hline
Bậc hai chuẩn & $\dfrac{K\omega_n^2}{s^2+2\xi\omega_ns+\omega_n^2}$ & phụ thuộc $\xi$ & pha tiến từ $0^\circ$ về gần $-180^\circ$\\
\hline
\end{tabular}
\renewcommand{\arraystretch}{1}
\end{center}

\subsection*{4. Khâu quán tính bậc nhất}

\[
\boxed{G(s)=\frac{K}{Ts+1}},\qquad T>0.
\]

Đáp ứng xung:
\[
\boxed{g(t)=\frac{K}{T}e^{-t/T}1(t)}.
\]

Đáp ứng bước:
\[
\boxed{h(t)=K(1-e^{-t/T})1(t)}.
\]

Tại $t=T$:
\[
h(T)=K(1-e^{-1})\approx0.632K.
\]
Vì vậy $T$ là \textbf{hằng số thời gian}.

Đặc tính tần số:
\[
G(\jj\omega)=\frac{K}{1+\jj\omega T},
\]
\[
\boxed{
M(\omega)=\frac{|K|}{\sqrt{1+(\omega T)^2}},
\qquad
\varphi(\omega)=-\tan^{-1}(\omega T)
}
\]
với $K>0$.

Tần số gãy:
\[
\boxed{\omega_b=\frac1T}.
\]
Sau tần số gãy, đường Bode biên độ tiệm cận có độ dốc $-20$ dB/dec.

\subsubsection*{Dạng bài: Nhận dạng khâu bậc nhất}

Đưa hàm truyền về:
\[
G(s)=\frac{K}{Ts+1}.
\]
Ví dụ:
\[
G(s)=\frac{6}{2s+4}
=\frac{3/2}{0.5s+1}.
\]
Suy ra:
\[
\boxed{K=\frac32,\qquad T=0.5}.
\]

\subsection*{5. Khâu bậc hai}

Dạng chuẩn:
\[
\boxed{
G(s)=
\frac{K\omega_n^2}
{s^2+2\xi\omega_ns+\omega_n^2}.
}
\]

Tương đương:
\[
G(s)=\frac{K}{T^2s^2+2\xi Ts+1},
\qquad
\omega_n=\frac1T.
\]

Cực:
\[
s_{1,2}
=
-\xi\omega_n
\pm
\omega_n\sqrt{\xi^2-1}.
\]

Nếu $0<\xi<1$:
\[
\omega_d=\omega_n\sqrt{1-\xi^2},
\]
\[
s_{1,2}=-\xi\omega_n\pm\jj\omega_d.
\]

Với $K=1$, $0<\xi<1$, đáp ứng xung:
\[
\boxed{
g(t)
=
\frac{\omega_n}{\sqrt{1-\xi^2}}
e^{-\xi\omega_nt}
\sin(\omega_dt)\,1(t).
}
\]

Một dạng đáp ứng bước tương đương:
\[
\boxed{
h(t)
=
1-
\frac{e^{-\xi\omega_nt}}{\sqrt{1-\xi^2}}
\sin(\omega_dt+\theta),
}
\]
trong đó:
\[
\theta=\cos^{-1}\xi.
\]

Ý nghĩa $\xi$:
\begin{itemize}[leftmargin=1.7em,itemsep=0.2em]
  \item $\xi=0$: dao động không tắt.
  \item $0<\xi<1$: dao động tắt dần.
  \item $\xi=1$: tắt dần tới hạn.
  \item $\xi>1$: quá tắt, hai cực thực âm phân biệt.
\end{itemize}

\subsubsection*{Dạng bài: Từ đa thức bậc hai tìm $\omega_n,\xi$}

Với:
\[
s^2+as+b,
\]
so sánh với:
\[
s^2+2\xi\omega_ns+\omega_n^2.
\]
Suy ra:
\[
\boxed{
\omega_n=\sqrt b,
\qquad
\xi=\frac{a}{2\sqrt b}.
}
\]

\subsection*{6. Khâu tích phân và vi phân trên Bode}

\textbf{Tích phân lý tưởng}
\[
G(s)=\frac1s
\Longrightarrow
G(\jj\omega)=\frac1{\jj\omega}.
\]
\[
\boxed{
L(\omega)=-20\log_{10}\omega,
\qquad
\varphi=-90^\circ.
}
\]
Độ dốc: $-20$ dB/dec trên toàn dải.

\textbf{Vi phân lý tưởng}
\[
G(s)=s
\Longrightarrow
G(\jj\omega)=\jj\omega.
\]
\[
\boxed{
L(\omega)=20\log_{10}\omega,
\qquad
\varphi=+90^\circ.
}
\]
Độ dốc: $+20$ dB/dec trên toàn dải.

\subsection*{7. Dạng bài Chương 4 và quy trình giải}

\subsubsection*{Dạng 1: Từ $G(s)$ tìm $g(t)$ và $h(t)$}

\[
\boxed{
g(t)=\Lap^{-1}\{G(s)\},
\qquad
h(t)=\Lap^{-1}\left\{\frac{G(s)}s\right\}.
}
\]
Nếu biểu thức hữu tỷ khó lấy ngược, phân tích phân thức như Chương 2.

\subsubsection*{Dạng 2: Từ $h(t)$ hoặc $g(t)$ suy ra $G(s)$}

\[
G(s)=\Lap\{g(t)\}.
\]
Nếu cho $h(t)$:
\[
H(s)=\Lap\{h(t)\},
\qquad
G(s)=sH(s)
\]
khi điều kiện đầu phù hợp.

\subsubsection*{Dạng 3: Tìm đặc tính tần số}

\begin{enumerate}[leftmargin=1.9em,itemsep=0.2em]
  \item Thay $s=\jj\omega$.
  \item Đưa về $P(\omega)+\jj Q(\omega)$ hoặc dạng tích.
  \item Tính $M(\omega)=|G(\jj\omega)|$.
  \item Tính pha bằng góc phần tư đúng.
  \item Nếu vẽ Bode: cộng biên độ dB và pha của từng khâu thành phần.
\end{enumerate}

% ============================================================
\clearpage
\section*{BẢNG TRA NHANH CHƯƠNG 1 -- 4}

\subsection*{A. Chuỗi thao tác theo dạng bài}

\begin{center}
\small
\renewcommand{\arraystretch}{1.35}
\begin{tabular}{|p{0.28\textwidth}|p{0.60\textwidth}|}
\hline
\textbf{Bài toán} & \textbf{Chuỗi thao tác}\\
\hline
PTVP $\to y(t)$ & Laplace hai vế $\to Y(s)$ $\to$ phân thức tối giản $\to$ Laplace ngược\\
\hline
PTVP $\to G(s)$ & điều kiện đầu bằng $0$ $\to$ Laplace $\to Y(s)/U(s)$\\
\hline
$G(s)\to g(t)$ & $g(t)=\Lap^{-1}\{G(s)\}$\\
\hline
$G(s)\to h(t)$ & $h(t)=\Lap^{-1}\{G(s)/s\}$\\
\hline
$h(t)\to G(s)$ & $H(s)=\Lap\{h(t)\}$, rồi $G(s)=sH(s)$\\
\hline
State-space $\to G(s)$ & $G(s)=C(sI-A)^{-1}B+D$\\
\hline
Sơ đồ khối $\to G_{eq}$ & vòng trong $\to$ nối tiếp / song song $\to$ hồi tiếp $\to$ lặp lại\\
\hline
Graph $\to G(s)$ & $P_k\to L_i\to\Delta\to\Delta_k\to$ Mason\\
\hline
$G(s)\to G(\jj\omega)$ & thay $s=\jj\omega\to M(\omega)\to\varphi(\omega)\to$ Bode\\
\hline
\end{tabular}
\renewcommand{\arraystretch}{1}
\end{center}

\subsection*{B. Công thức bắt buộc thuộc}

\[
\boxed{\Lap\{f'(t)\}=sF(s)-f(0^+)}
\]
\[
\boxed{\Lap\{f''(t)\}=s^2F(s)-sf(0^+)-f'(0^+)}
\]
\[
\boxed{G(s)=\frac{Y(s)}{U(s)}\quad\text{(điều kiện đầu bằng 0)}}
\]
\[
\boxed{G_{cl}(s)=\frac{G(s)}{1+G(s)H(s)}\quad\text{(hồi tiếp âm)}}
\]
\[
\boxed{G(s)=C(sI-A)^{-1}B+D}
\]
\[
\boxed{g(t)=\Lap^{-1}\{G(s)\},
\qquad h(t)=\Lap^{-1}\{G(s)/s\}}
\]
\[
\boxed{G(\jj\omega)=G(s)\big|_{s=\jj\omega}}
\]
\[
\boxed{
\frac{C(s)}{R(s)}
=
\frac{\sum_kP_k\Delta_k}{\Delta}
\quad\text{(Mason)}.
}
\]

\subsection*{C. Checklist trước khi chốt đáp án}

\begin{itemize}[leftmargin=1.7em,itemsep=0.2em]
  \item Đã xác định đúng tín hiệu vào, đầu ra và dấu hồi tiếp chưa?
  \item Khi tìm hàm truyền, đã cho điều kiện đầu bằng $0$ chưa?
  \item Khi giải PTVP, đã thay đầy đủ $y(0),y'(0),\ldots$ trong công thức Laplace chưa?
  \item Phân thức Laplace ngược đã hợp thức chưa? Nếu chưa, đã chia đa thức chưa?
  \item Mạch điện đã dùng đúng $Z_R=R$, $Z_L=sL$, $Z_C=1/(sC)$ khi điều kiện đầu bằng $0$ chưa?
  \item Sơ đồ khối đã rút vòng trong trước vòng ngoài chưa?
  \item Mason đã liệt kê đủ đường truyền thuận, vòng lặp và các vòng không chạm chưa?
  \item Khi tìm pha của $G(\jj\omega)$, đã xét đúng góc phần tư chưa?
  \item Kết quả cuối đã có đúng đơn vị, điều kiện và ký hiệu chưa?
\end{itemize}

\medskip
\noindent\textit{Ghi chú tích hợp: file này được cố ý để độc lập với phần nội dung chính.
Khi hoàn thiện main.tex, chỉ cần đặt lệnh
\texttt{\string\input\{tong\_hop\_chuong\_1\_4.tex\}}
ngay trước \texttt{\string\end\{document\}}.}
